% GENERATED FILE. Edit canonical lessons, then run npm run generate:books.
% canonical-source-hash: fnv1a64:0ae94a65f277d416
% canonical-lessons: RU-C281-korrespondent, RU-C281-reportazh, RU-C281-vedushchiy, RU-C281-efir, RU-C281-telekanal

\chapter{On the air}
\label{ch:ru-on-the-air}
\hlchaptermodality{281}

\begin{chapteropening}
\textbf{By the end of this chapter:} \emph{I can talk about a correspondent, a news report, a presenter, being on air and a TV channel.}

Complete the last lesson of chapter 281: I can talk about a correspondent, a news report, a presenter, being on air and a TV channel.
\end{chapteropening}

\section[\texorpdfstring{\ru{корреспондент} (korrespondént)}{корреспондент (korrespondént)}]{\ru{корреспондент} (korrespondént) — a correspondent}

\label{lesson:RU-C281-korrespondent}



\begin{quote}
Before the new one: say the Russian for to confirm, then the Russian for to refute.
\end{quote}

\subsection*{You'll want to know: \ru{корреспондент}}
\textbf{\ru{корреспондент}} (\emph{korrespondént}) — \textquotedblleft{}a correspondent\textquotedblright{}.

\begin{quote}
\ru{Корреспондент} \ru{поехал} \ru{на} \ru{север}.
\end{quote}

\begin{quote}
\emph{Korrespondént poyékhal na séver.}
\end{quote}

\begin{quote}
The correspondent went north.
\end{quote}

Also \emph{\ru{журналист}}. Feminine forms are colloquial, and \emph{\ru{корреспондент}} serves for a woman too.

\begin{cousinweb}
\emph{\ru{корреспондент}} is usually traced through German \emph{Korrespondent} or French \emph{correspondant} to Medieval Latin \emph{correspondēre} \textquotedblleft{}to answer together\textquotedblright{}; EtymDB lists no parent.
\end{cousinweb}

\begin{grammarlens}[title={What you've built}]
One more word for telling, explaining and describing at B1.
\end{grammarlens}

\subsection*{Your turn}
\begin{itemize}
\raggedright
  \item \emph{Say it:} \emph{korrespondént}
  \item \emph{Say it:} \emph{korrespondént}, once more
  \item \emph{Say it:} a sentence of your own with \emph{korrespondént}
  \item \emph{Recall:} say the Russian for to confirm, then the Russian for to refute, then say \emph{korrespondént} again
\end{itemize}

\begin{tcolorbox}[breakable,colback=teal!4,colframe=teal!35!black,title={Before you move on}]
What does \ru{корреспондент} mean? (\textquotedblleft{}A correspondent\textquotedblright{}.) Say it once more.
\end{tcolorbox}
\section[\texorpdfstring{\ru{репортаж} (reportázh)}{репортаж (reportázh)}]{\ru{репортаж} (reportázh) — a report (on TV or radio), a news story}

\label{lesson:RU-C281-reportazh}



\begin{quote}
Before the new one: say the Russian for to refute, then the Russian for a correspondent.
\end{quote}

\subsection*{You'll want to know: \ru{репортаж}}
\textbf{\ru{репортаж}} (\emph{reportázh}) — \textquotedblleft{}a report (on TV or radio), a news story\textquotedblright{}.

\begin{quote}
\ru{Вечером} \ru{мы} \ru{смотрели} \ru{репортаж} \ru{о} \ru{футболе}.
\end{quote}

\begin{quote}
\emph{Vécherom my smotréli reportázh o futbóle.}
\end{quote}

\begin{quote}
In the evening we watched a report about football.
\end{quote}

Masculine. \emph{\ru{Вести} \ru{репортаж}} — to report live.

\begin{cousinweb}
\emph{\ru{репортаж}} was borrowed from French \emph{reportage} (EtymDB), built on \emph{reporter} \textquotedblleft{}to bring back, to report\textquotedblright{} and the ending \emph{-age}; the verb goes back to Latin \emph{reportāre}, from \emph{portāre} \textquotedblleft{}to carry\textquotedblright{}.
\end{cousinweb}

\begin{grammarlens}[title={What you've built}]
One more word for telling, explaining and describing at B1.
\end{grammarlens}

\subsection*{Your turn}
\begin{itemize}
\raggedright
  \item \emph{Say it:} \emph{reportázh}
  \item \emph{Say it:} \emph{reportázh}, once more
  \item \emph{Say it:} a sentence of your own with \emph{reportázh}
  \item \emph{Recall:} say the Russian for to refute, then the Russian for a correspondent, then say \emph{reportázh} again
\end{itemize}

\begin{tcolorbox}[breakable,colback=teal!4,colframe=teal!35!black,title={Before you move on}]
What does \ru{репортаж} mean? (\textquotedblleft{}A report (on TV or radio), a news story\textquotedblright{}.) Say it once more.
\end{tcolorbox}
\section[\texorpdfstring{\ru{ведущий} (vedúshchiy)}{ведущий (vedúshchiy)}]{\ru{ведущий} (vedúshchiy) — a presenter, a host}

\label{lesson:RU-C281-vedushchiy}



\begin{quote}
Before the new one: say the Russian for a correspondent, then the Russian for a report.
\end{quote}

\subsection*{You'll want to know: \ru{ведущий}}
\textbf{\ru{ведущий}} (\emph{vedúshchiy}) — \textquotedblleft{}a presenter, a host\textquotedblright{}.

\begin{quote}
\ru{Ведущий} \ru{программы} \ru{очень} \ru{весёлый}.
\end{quote}

\begin{quote}
\emph{Vedúshchiy prográmmy óchen vesyólyy.}
\end{quote}

\begin{quote}
The presenter of the programme is very cheerful.
\end{quote}

A participle of \emph{\ru{вести}} \textquotedblleft{}to lead\textquotedblright{} used as a noun; a woman is \emph{\ru{ведущая}}. It declines like an adjective.

\begin{cousinweb}
\emph{\ru{ведущий}} is the present participle of \emph{\ru{вести}} \textquotedblleft{}to lead\textquotedblright{} (EtymDB), an inherited Slavic verb from the same root as \emph{\ru{водить}}: the one who leads the show.
\end{cousinweb}

\begin{grammarlens}[title={What you've built}]
One more word for telling, explaining and describing at B1.
\end{grammarlens}

\subsection*{Your turn}
\begin{itemize}
\raggedright
  \item \emph{Say it:} \emph{vedúshchiy}
  \item \emph{Say it:} \emph{vedúshchiy}, once more
  \item \emph{Say it:} a sentence of your own with \emph{vedúshchiy}
  \item \emph{Recall:} say the Russian for a correspondent, then the Russian for a report, then say \emph{vedúshchiy} again
\end{itemize}

\begin{tcolorbox}[breakable,colback=teal!4,colframe=teal!35!black,title={Before you move on}]
What does \ru{ведущий} mean? (\textquotedblleft{}A presenter, a host\textquotedblright{}.) Say it once more.
\end{tcolorbox}
\section[\texorpdfstring{\ru{эфир} (efír)}{эфир (efír)}]{\ru{эфир} (efír) — the air (broadcasting)}

\label{lesson:RU-C281-efir}



\begin{quote}
Before the new one: say the Russian for a report, then the Russian for a presenter.
\end{quote}

\subsection*{You'll want to know: \ru{эфир}}
\textbf{\ru{эфир}} (\emph{efír}) — \textquotedblleft{}the air (broadcasting)\textquotedblright{}.

\begin{quote}
\ru{Программа} \ru{вышла} \ru{в} \ru{эфир} \ru{в} \ru{семь} \ru{часов}.
\end{quote}

\begin{quote}
\emph{Prográmma výshla v efír v sem chasóv.}
\end{quote}

\begin{quote}
The programme went on air at seven.
\end{quote}

\emph{\ru{В} \ru{прямом} \ru{эфире}} — live. In chemistry it also means ether.

\begin{cousinweb}
\emph{\ru{эфир}} was borrowed from Greek \emph{aith\'{\={e}}r} \textquotedblleft{}the upper air\textquotedblright{} (EtymDB), built on \emph{aíthō} \textquotedblleft{}to burn, to shine\textquotedblright{}.
\end{cousinweb}

\begin{grammarlens}[title={What you've built}]
One more word for telling, explaining and describing at B1.
\end{grammarlens}

\subsection*{Your turn}
\begin{itemize}
\raggedright
  \item \emph{Say it:} \emph{efír}
  \item \emph{Say it:} \emph{efír}, once more
  \item \emph{Say it:} a sentence of your own with \emph{efír}
  \item \emph{Recall:} say the Russian for a report, then the Russian for a presenter, then say \emph{efír} again
\end{itemize}

\begin{tcolorbox}[breakable,colback=teal!4,colframe=teal!35!black,title={Before you move on}]
What does \ru{эфир} mean? (\textquotedblleft{}The air (broadcasting)\textquotedblright{}.) Say it once more.
\end{tcolorbox}
\section[\texorpdfstring{\ru{телеканал} (telekanál)}{телеканал (telekanál)}]{\ru{телеканал} (telekanál) — a TV channel}

\label{lesson:RU-C281-telekanal}



\begin{quote}
Before the new one: say the Russian for a presenter, then the Russian for the air.
\end{quote}

\subsection*{You'll want to know: \ru{телеканал}}
\textbf{\ru{телеканал}} (\emph{telekanál}) — \textquotedblleft{}a TV channel\textquotedblright{}.

\begin{quote}
\ru{Этот} \ru{телеканал} \ru{показывает} \ru{много} \ru{фильмов}.
\end{quote}

\begin{quote}
\emph{Étot telekanál pokázyvayet mnógo fílmov.}
\end{quote}

\begin{quote}
This TV channel shows a lot of films.
\end{quote}

Masculine. Short for \emph{\ru{телевизионный} \ru{канал}}.

\begin{cousinweb}
\emph{\ru{телеканал}} joins \emph{\ru{теле}-}, from Greek \emph{tēle} \textquotedblleft{}far off\textquotedblright{}, and \emph{\ru{канал}} \textquotedblleft{}a channel\textquotedblright{}, which comes from Latin \emph{canālis} \textquotedblleft{}a pipe, a channel\textquotedblright{} (EtymDB).
\end{cousinweb}

\begin{grammarlens}[title={What you've built}]
One more word for telling, explaining and describing at B1.
\end{grammarlens}

\subsection*{Your turn}
\begin{itemize}
\raggedright
  \item \emph{Say it:} \emph{telekanál}
  \item \emph{Say it:} \emph{telekanál}, once more
  \item \emph{Say it:} a sentence of your own with \emph{telekanál}
  \item \emph{Recall:} say the Russian for a presenter, then the Russian for the air, then say \emph{telekanál} again
\end{itemize}

\begin{tcolorbox}[breakable,colback=teal!4,colframe=teal!35!black,title={Before you move on}]
What does \ru{телеканал} mean? (\textquotedblleft{}A TV channel\textquotedblright{}.) Say it once more.
\end{tcolorbox}
